\documentclass[a4paper,11pt]{article}

% =========================================================================
% MAKHASWA HOLDINGS | BULK EARTHWORKS (TWO-PAGE STANDALONE SECTION)
% Compile with pdfLaTeX on Overleaf. Upload the three referenced photos and
% logo.jpeg to the project's images/ folder if you want the photographs shown.
% The document also compiles with labelled placeholders if photos are missing.
% =========================================================================

\usepackage[utf8]{inputenc}
\usepackage[T1]{fontenc}
\usepackage[left=1.2cm,right=1.2cm,top=0cm,bottom=1.4cm,footskip=0.9cm]{geometry}
\usepackage{helvet}
\renewcommand{\familydefault}{\sfdefault}
\usepackage{graphicx}
\usepackage{xcolor}
\usepackage{tikz}
\usetikzlibrary{calc,shadows}
\usepackage[most]{tcolorbox}
\usepackage{tabularx}
\usepackage{colortbl}
\usepackage{array}
\usepackage{fancyhdr}
\usepackage{hyperref}
\usepackage{fontawesome5}
\providecommand{\faShieldAlt}{\faShield*}
\graphicspath{{images/}{./}}

% Corporate colour palette (matched to the company profile)
\definecolor{mhnavy}{HTML}{082D5A}
\definecolor{mhnavy2}{HTML}{0C3A6D}
\definecolor{mhnavy3}{HTML}{124884}
\definecolor{mhblue}{HTML}{1A56A0}
\definecolor{mhbluebright}{HTML}{0066CC}
\definecolor{mhbluelight}{HTML}{F0F4FA}
\definecolor{mhgold}{HTML}{C8A951}
\definecolor{mhoffwhite}{HTML}{F8FAFC}
\definecolor{mhlightgray}{HTML}{E2E8F0}
\definecolor{mhtext}{HTML}{1E293B}
\definecolor{mhtextmuted}{HTML}{64748B}
\pagecolor{white}
\color{mhtext}

\hypersetup{colorlinks=true,linkcolor=mhnavy,urlcolor=mhblue,
    pdftitle={Makhaswa Holdings - Bulk Earthworks},
    pdfauthor={Makhaswa Holdings (Pty) Ltd}}

% Section header: same navy band, gold rule, logo treatment
\newcommand{\sectionbanner}[2]{%
    \begin{tikzpicture}[remember picture,overlay]
        \fill[mhnavy] (current page.north west) rectangle
            ([yshift=-3.5cm]current page.north east);
        \fill[mhgold] ([yshift=-3.5cm]current page.north west) rectangle
            ([yshift=-3.38cm]current page.north east);
        \node[anchor=west,inner sep=0pt,text width=0.64\paperwidth]
            at ([xshift=1.2cm,yshift=-1.75cm]current page.north west) {%
                {\fontsize{21pt}{25pt}\selectfont\bfseries\color{white}#1}\\[0.22cm]
                {\footnotesize\bfseries\color{mhbluelight}\MakeUppercase{#2}}
            };
        \node[anchor=east,inner sep=0pt]
            at ([xshift=-1.2cm,yshift=-1.75cm]current page.north east) {%
                \begin{tcolorbox}[enhanced,colback=white,colframe=mhlightgray,
                    boxrule=0.8pt,arc=3pt,boxsep=2pt,
                    left=4pt,right=4pt,top=3pt,bottom=3pt,
                    width=3.2cm,halign=center,drop shadow=black!15]
                    \IfFileExists{images/logo.jpeg}{%
                        \includegraphics[height=1.0cm,keepaspectratio]{images/logo.jpeg}%
                    }{%
                        {\fontsize{11pt}{12pt}\selectfont\bfseries\color{mhnavy}MAKHASWA}\\[-1pt]
                        {\fontsize{7pt}{8pt}\selectfont\color{mhblue}HOLDINGS}%
                    }
                \end{tcolorbox}
            };
    \end{tikzpicture}%
    \vspace*{3.65cm}\noindent
}

% Image helper, with a readable placeholder for missing image assets
\newcommand{\earthimage}[2]{%
    \IfFileExists{images/#1}{%
        \includegraphics[width=\linewidth,height=#2,keepaspectratio=false]{images/#1}%
    }{%
        \fcolorbox{mhlightgray}{mhbluelight}{%
            \parbox[c][#2][c]{\dimexpr\linewidth-2\fboxsep-2\fboxrule\relax}{%
                \centering\fontsize{7pt}{9pt}\selectfont\color{mhtextmuted}
                \faImage\quad Add photo:\\[2pt]#1}}
    }%
}

\newcommand{\servicecard}[3][\faCogs]{%
    \begin{tcolorbox}[
        enhanced,colback=mhoffwhite,colframe=mhlightgray,boxrule=0.8pt,
        arc=4pt,left=8pt,right=8pt,top=8pt,bottom=8pt,
        borderline west={3pt}{0pt}{mhblue},drop shadow=black!3]
        {\fontsize{9.8pt}{12pt}\selectfont\bfseries\color{mhnavy}%
        {\color{mhblue}#1}\quad #2}\\[3pt]
        {\fontsize{8.4pt}{11.6pt}\selectfont\color{mhtext}#3}
    \end{tcolorbox}%
}

\newcommand{\processcard}[4]{%
    \begin{tcolorbox}[enhanced,colback=white,colframe=mhlightgray,
        boxrule=0.8pt,arc=4pt,left=3pt,right=3pt,top=6pt,bottom=6pt,
        halign=center,borderline north={2.4pt}{0pt}{#1},
        drop shadow=black!4]
        {\fontsize{11pt}{13pt}\selectfont\color{#1}#4}\\[2pt]
        {\fontsize{12.5pt}{15pt}\selectfont\bfseries\color{#1}#2}\\[2pt]
        {\fontsize{7pt}{8.5pt}\selectfont\color{mhtextmuted}\MakeUppercase{#3}}
    \end{tcolorbox}%
}

\pagestyle{fancy}
\fancyhf{}
\renewcommand{\headrulewidth}{0pt}
\fancyfoot[L]{%
    \fontsize{8.2pt}{10.5pt}\selectfont\color{mhtextmuted}%
    \textbf{Makhaswa Holdings (Pty) Ltd} $\bullet$
    CIDB Grade 8 CE PE / 8 ME PE $\bullet$ Reg: K2017/307900/07}
\fancyfoot[R]{%
    \fontsize{8.2pt}{10.5pt}\selectfont\bfseries\color{mhnavy}\thepage}

\begin{document}

% =========================================================================
% PAGE 1 --- BULK EARTHWORKS CAPABILITY
% =========================================================================
\sectionbanner{BULK EARTHWORKS}{Mass Excavation, Cut-to-Fill Balancing \& Site Preparation}

\noindent
\begin{minipage}[t]{0.62\textwidth}
    \vspace{0pt}
    {\large\bfseries\color{mhnavy} Bulk Earthworks \& Ground Engineering Capability}
    \vspace{0.12cm}

    {\fontsize{9pt}{13pt}\selectfont\color{mhtext}
    \textbf{Makhaswa Holdings (Pty) Ltd} undertakes bulk earthworks for civil
    infrastructure, road construction, commercial developments and general
    building projects. The work includes the excavation, movement, placement
    and shaping of soil and rock to prepare stable construction platforms,
    foundations, road formations and service corridors.
    }

    \vspace{0.22cm}

    {\fontsize{9pt}{13pt}\selectfont\color{mhtext}
    Our scope covers mass excavation, cut-to-fill balancing, site clearance,
    topsoil stripping, soil compaction, ground stabilisation, stockpiling and
    haulage. Excavated material is assessed for suitable reuse in engineered
    fill, with surplus and unsuitable material managed in accordance with
    project requirements.
    }

    \vspace{0.22cm}

    {\fontsize{9pt}{13pt}\selectfont\color{mhtext}
    Earthworks are integrated with road layerworks, stormwater drainage,
    sewer pipelines and building foundations. Survey levels, fill material,
    layer thickness and compaction requirements are controlled against the
    engineer's drawings and specifications.
    }
\end{minipage}
\hfill
\begin{minipage}[t]{0.34\textwidth}
    \vspace{0pt}
    \begin{tcolorbox}[
        enhanced,colback=mhoffwhite,colframe=mhnavy,boxrule=1pt,arc=4pt,
        title={\bfseries\color{white}\faTractor\ Earthworks Scope of Work},
        colbacktitle=mhnavy,left=7pt,right=7pt,top=6pt,bottom=6pt,
        drop shadow=black!6]
        \fontsize{8.1pt}{11.6pt}\selectfont
        \textcolor{mhnavy}{\rule{4pt}{4pt}}\hspace{4pt}\textbf{Mass excavation}\\
        \textcolor{mhnavy}{\rule{4pt}{4pt}}\hspace{4pt}\textbf{Cut-to-fill balancing}\\
        \textcolor{mhnavy}{\rule{4pt}{4pt}}\hspace{4pt}\textbf{Site clearance \& stripping}\\
        \textcolor{mhnavy}{\rule{4pt}{4pt}}\hspace{4pt}\textbf{Compaction \& stabilisation}\\
        \textcolor{mhnavy}{\rule{4pt}{4pt}}\hspace{4pt}\textbf{Bulk platform preparation}\\
        \textcolor{mhnavy}{\rule{4pt}{4pt}}\hspace{4pt}\textbf{Stockpiling \& haulage}\\
        \textcolor{mhnavy}{\rule{4pt}{4pt}}\hspace{4pt}\textbf{Controlled demolition}
    \end{tcolorbox}

    \vspace{0.20cm}

    \begin{tcolorbox}[
        enhanced,colback=white,colframe=mhlightgray,boxrule=0.8pt,arc=4pt,
        left=6pt,right=6pt,top=6pt,bottom=6pt,halign=center,
        drop shadow=black!4]
        {\fontsize{18pt}{20pt}\selectfont\bfseries\color{mhblue}8 CE PE}\\[1pt]
        {\fontsize{7.4pt}{9pt}\selectfont\color{mhtextmuted}%
        \MakeUppercase{CIDB Civil Engineering Grade}}\\[4pt]
        \textcolor{mhlightgray}{\rule{3.2cm}{0.6pt}}\\[4pt]
        {\fontsize{13pt}{15pt}\selectfont\bfseries\color{mhnavy}Mod AASHTO}\\[1pt]
        {\fontsize{7.2pt}{9pt}\selectfont\color{mhtextmuted}%
        \MakeUppercase{Compaction Test Reference}}
    \end{tcolorbox}
\end{minipage}

\vfill

% Photos already referenced in the source company profile
\noindent
\begin{minipage}[t]{0.315\textwidth}
    \begin{tcolorbox}[enhanced,colback=white,colframe=mhlightgray,
        boxrule=0.8pt,arc=4pt,left=0pt,right=0pt,top=0pt,bottom=5pt,
        drop shadow=black!4]
        \earthimage{bulk-trenching-earthworks-1.jpeg}{2.8cm}
        \centering\vspace{3pt}
        {\fontsize{8pt}{9.8pt}\selectfont\bfseries\color{mhnavy}%
        Mass Excavation \& Trenching}\\[1pt]
        {\fontsize{7pt}{8.4pt}\selectfont\color{mhtextmuted}%
        Excavation and Subgrade Preparation}
    \end{tcolorbox}
\end{minipage}
\hfill
\begin{minipage}[t]{0.315\textwidth}
    \begin{tcolorbox}[enhanced,colback=white,colframe=mhlightgray,
        boxrule=0.8pt,arc=4pt,left=0pt,right=0pt,top=0pt,bottom=5pt,
        drop shadow=black!4]
        \earthimage{plant-and-equipment-fleet.jpeg}{2.8cm}
        \centering\vspace{3pt}
        {\fontsize{8pt}{9.8pt}\selectfont\bfseries\color{mhnavy}%
        Earthmoving Plant}\\[1pt]
        {\fontsize{7pt}{8.4pt}\selectfont\color{mhtextmuted}%
        Heavy Equipment and Material Handling}
    \end{tcolorbox}
\end{minipage}
\hfill
\begin{minipage}[t]{0.315\textwidth}
    \begin{tcolorbox}[enhanced,colback=white,colframe=mhlightgray,
        boxrule=0.8pt,arc=4pt,left=0pt,right=0pt,top=0pt,bottom=5pt,
        drop shadow=black!4]
        \earthimage{earthworks-surfacing-watertruck.jpeg}{2.8cm}
        \centering\vspace{3pt}
        {\fontsize{8pt}{9.8pt}\selectfont\bfseries\color{mhnavy}%
        Earthworks \& Moisture Control}\\[1pt]
        {\fontsize{7pt}{8.4pt}\selectfont\color{mhtextmuted}%
        Water Truck and Site Operations}
    \end{tcolorbox}
\end{minipage}

\vfill

\noindent
{\large\bfseries\color{mhnavy} Core Earthworks Services}
\vspace{0.10cm}

\noindent
\begin{minipage}[t]{0.485\textwidth}
    \servicecard[\faTractor]{01. Mass Excavation}{Bulk removal of soil and rock for
    foundation excavations, basement levels, road cuttings and building
    platforms, to the required lines and levels.}
    \vspace{0.16cm}
    \servicecard[\faBalanceScale]{02. Cut-to-Fill Balancing}{Calculation and
    distribution of suitable excavated material between cut and fill areas
    to reduce unnecessary off-site disposal and imported fill.}
    \vspace{0.16cm}
    \servicecard[\faLeaf]{03. Site Clearance \& Topsoil Stripping}{Removal of
    vegetation, obstructions and unsuitable surface material. Reusable
    topsoil is separately stockpiled for landscaping or reinstatement.}
\end{minipage}
\hfill
\begin{minipage}[t]{0.485\textwidth}
    \servicecard[\faLayerGroup]{04. Compaction \& Soil Stabilisation}{Placement
    and compaction of engineered fill in specified layers; lime or cement
    stabilisation where required by the approved material design.}
    \vspace{0.16cm}
    \servicecard[\faBuilding]{05. Bulk Platform \& Formation Preparation}{Formation
    of level terraces, industrial and commercial building footprints,
    road subgrades and earthworks for infrastructure services.}
    \vspace{0.16cm}
    \servicecard[\faTruckMoving]{06. Stockpiling \& Haulage}{Segregation,
    loading, transport and management of reusable, excess and unsuitable
    material, including removal to authorised disposal facilities.}
\end{minipage}

\clearpage

% =========================================================================
% PAGE 2 --- EXECUTION METHOD, TECHNICAL CONTROLS & ASSOCIATED CIVILS
% =========================================================================
\sectionbanner{EARTHWORKS EXECUTION}{Construction Sequence, Plant Operations \& Quality Control}

\noindent
{\large\bfseries\color{mhnavy} Bulk Earthworks Construction Sequence}
\vspace{0.12cm}

\noindent
\begin{minipage}[t]{0.185\textwidth}
    \processcard{mhnavy}{Survey}{Set-out \& levels}{\faRulerCombined}
\end{minipage}\hfill
\begin{minipage}[t]{0.185\textwidth}
    \processcard{mhblue}{Clear}{Strip \& stockpile}{\faLeaf}
\end{minipage}\hfill
\begin{minipage}[t]{0.185\textwidth}
    \processcard{mhbluebright}{Excavate}{Cut \& load}{\faTractor}
\end{minipage}\hfill
\begin{minipage}[t]{0.185\textwidth}
    \processcard{mhnavy3}{Place}{Fill \& grade}{\faLayerGroup}
\end{minipage}\hfill
\begin{minipage}[t]{0.185\textwidth}
    \processcard{mhblue}{Test}{Compact \& verify}{\faCheckCircle}
\end{minipage}

\vspace{0.20cm}

\noindent
\begin{tcolorbox}[enhanced,colback=mhoffwhite,colframe=mhlightgray,
    boxrule=0.7pt,arc=4pt,left=9pt,right=9pt,top=7pt,bottom=7pt,
    borderline west={3pt}{0pt}{mhgold},drop shadow=black!4]
    {\fontsize{8.1pt}{11.0pt}\selectfont\color{mhtext}
    \textbf{Execution method:} Survey control and site clearance precede
    excavation. Suitable material is directed to designated fill areas;
    unsuitable material is separated. Fill is spread, moisture-conditioned
    where necessary, compacted in specified layers and checked for line,
    level and required density before acceptance of the next layer.
    }
\end{tcolorbox}

\vfill

\noindent
{\large\bfseries\color{mhnavy} Earthworks Activities \& Technical Controls}
\vspace{0.10cm}

{\fontsize{7.7pt}{9.8pt}\selectfont
\renewcommand{\arraystretch}{1.23}
\noindent
\begin{tabularx}{\textwidth}{@{}p{3.3cm} p{6.6cm} X@{}}
\rowcolor{mhnavy}
{\color{white}\bfseries Work Item} &
{\color{white}\bfseries Construction Method} &
{\color{white}\bfseries Site Control / Inspection} \\
\rowcolor{mhbluelight}
\bfseries Survey \& setting out & Establish existing ground levels, cut/fill boundaries, excavation limits and finished platform levels & Survey records, benchmarks and level verification \\
\bfseries Clearing \& stripping & Remove vegetation, topsoil and obstructions; segregate and stockpile recoverable topsoil & Approved stripping limits and stockpile locations \\
\rowcolor{mhbluelight}
\bfseries Excavation \& material selection & Excavate soil or rock; assess material suitability for fill and direct loads to designated areas & Excavation profiles, material classification and approved disposal routes \\
\bfseries Fill placement \& compaction & Place approved fill in controlled layers, moisture-condition and compact using appropriate rollers & Layer thickness, density testing and engineer's specified compaction criteria \\
\rowcolor{mhbluelight}
\bfseries Stabilisation \& subgrade & Rip, shape and treat unsuitable ground with lime or cement where prescribed & Approved mix design, treatment depth, levels and density testing \\
\bfseries Controlled demolition & Remove existing structures and concrete obstructions under an approved demolition method; separate reusable rubble where specified & Demolition method, service isolation, waste segregation and disposal documentation \\
\end{tabularx}
\renewcommand{\arraystretch}{1}
}

\vfill

\noindent
\begin{minipage}[t]{0.485\textwidth}
    \begin{tcolorbox}[enhanced,colback=white,colframe=mhlightgray,
        boxrule=0.8pt,arc=4pt,
        title={\bfseries\color{white}\faTruckMonster\ Plant \& Material Logistics},
        colbacktitle=mhnavy,left=8pt,right=8pt,top=7pt,bottom=7pt,
        drop shadow=black!4]
        {\fontsize{8.1pt}{11.3pt}\selectfont\color{mhtext}
        Excavators, bulldozers, graders, tipper trucks, water trucks and
        compaction rollers are selected to suit ground conditions and haul
        distances. Excavation, stockpiling, loading, haulage and material
        placement are coordinated to maintain the approved work sequence.
        }
    \end{tcolorbox}
\end{minipage}
\hfill
\begin{minipage}[t]{0.485\textwidth}
    \begin{tcolorbox}[enhanced,colback=white,colframe=mhlightgray,
        boxrule=0.8pt,arc=4pt,
        title={\bfseries\color{white}\faClipboardCheck\ Quality \& Survey Control},
        colbacktitle=mhblue,left=8pt,right=8pt,top=7pt,bottom=7pt,
        drop shadow=black!4]
        {\fontsize{8.1pt}{11.3pt}\selectfont\color{mhtext}
        Setting out, as-built levels, material classification and layer
        compaction are checked against the project specification.
        In-situ density testing is assessed against the specified
        \textbf{Modified AASHTO} requirement, with test results recorded
        in the Inspection and Test Plan.
        }
    \end{tcolorbox}
\end{minipage}

\vspace{0.18cm}

\noindent
\begin{minipage}[t]{0.485\textwidth}
    \begin{tcolorbox}[enhanced,colback=white,colframe=mhlightgray,
        boxrule=0.8pt,arc=4pt,
        title={\bfseries\color{white}\faShieldAlt\ Safety \& Environmental Controls},
        colbacktitle=mhnavy,left=8pt,right=8pt,top=7pt,bottom=7pt,
        drop shadow=black!4]
        {\fontsize{8.1pt}{11.3pt}\selectfont\color{mhtext}
        Excavation stability, plant movement, buried services, dust,
        drainage and erosion are managed through the site risk assessment.
        Topsoil preservation, runoff protection and lawful disposal of
        waste material form part of the work method.
        }
    \end{tcolorbox}
\end{minipage}
\hfill
\begin{minipage}[t]{0.485\textwidth}
    \begin{tcolorbox}[enhanced,colback=white,colframe=mhlightgray,
        boxrule=0.8pt,arc=4pt,
        title={\bfseries\color{white}\faRoad\ Related Civil Engineering Works},
        colbacktitle=mhblue,left=8pt,right=8pt,top=7pt,bottom=7pt,
        drop shadow=black!4]
        {\fontsize{8.1pt}{11.3pt}\selectfont\color{mhtext}
        Earthworks support road construction, stormwater channels,
        drainage structures, sewer and bulk service trenches, foundation
        excavations and general building platforms. Work is sequenced
        with project management and the relevant civil trades.
        }
    \end{tcolorbox}
\end{minipage}

\vfill

\noindent
\begin{tcolorbox}[enhanced,colback=white,colframe=mhnavy,
    boxrule=0.8pt,arc=3pt,left=8pt,right=8pt,top=6pt,bottom=6pt,
    drop shadow=black!3]
    {\fontsize{8pt}{10.6pt}\selectfont\color{mhnavy}\bfseries
    \faCheckCircle\ Technical requirement:}\quad
    {\fontsize{8pt}{10.6pt}\selectfont\color{mhtext}
    Excavation depths, benching, layer thicknesses, material suitability,
    stabilisation rates and compaction targets are set by the engineer's
    drawings, geotechnical recommendations and project specification.
    }
\end{tcolorbox}

\end{document}
